\documentclass[12pt]{article}
\usepackage[utf8]{inputenc}
\usepackage[T1]{fontenc}
\usepackage[margin=1in]{geometry}
\usepackage{setspace}
\usepackage{xcolor}
\usepackage[colorlinks=true, urlcolor=blue!60!black, linkcolor=blue!60!black]{hyperref}
\setlength{\parindent}{2em}
\setlength{\parskip}{1.25em}
\renewcommand{\baselinestretch}{1.0}
\begin{document}
\begin{center}
{\LARGE \textit{Ziglar v. Abbasi}, 582 U.S. 120 (2017)\par}
\vspace{0.8em}
{\large Case Brief --- Four-Page Edition\par}
\end{center}
\vspace{1em}

\section{Memory Jogger}

Post-9/11 immigration detainees sued the Attorney General, FBI Director, and jail wardens for damages; the Court held 4--2 that extending \textit{Bivens} is disfavored, that claims in a new context fail if special factors counsel hesitation, and that no \textit{Bivens} remedy lay for challenges to high-level detention policy (\href{https://www.law.cornell.edu/supremecourt/text/15-1358}{Cornell LII}).

\section{Detailed Case Facts}

After the September 11 attacks, the FBI arrested over 700 noncitizens on immigration violations, holding those ``of interest'' without bail under a hold-until-cleared policy. Respondents, six men of Arab or South Asian descent, were among about 84 detainees held in the Administrative Maximum Special Housing Unit of Brooklyn's Metropolitan Detention Center for three to eight months. They alleged that official policy kept them in tiny cells over 23 hours a day with lights on around the clock, denied hygiene items and communication, and required frequent strip searches, and that wardens knowingly tolerated guards' abuse, including slamming detainees into walls and religious taunts.

\section{Procedural History}

Respondents sued former Attorney General Ashcroft, FBI Director Mueller, and INS Commissioner Ziglar (the Executive Officials) and Wardens Hasty and Sherman, asserting \textit{Bivens} claims challenging the detention policy, strip searches, and abuse, plus a conspiracy claim under 42 U.S.C. \S~1985(3). After \textit{Ashcroft v. Iqbal}, the district court dismissed the claims against the Executive Officials. A divided Second Circuit reinstated them and largely affirmed as to the Wardens.

\section{Judicial Votes}

Reversed in part, vacated and remanded in part, \textbf{4--2} (decided June 19, 2017).
\begin{itemize}
\item \textbf{Majority:} Justice Kennedy, joined by Chief Justice Roberts and Justices Thomas and Alito (Part IV--B a plurality without Thomas).
\item \textbf{Concurring in part and in the judgment:} Justice Thomas.
\item \textbf{Dissent:} Justice Breyer, joined by Justice Ginsburg.
\item \textbf{Not participating:} Justices Sotomayor, Kagan, and Gorsuch.
\end{itemize}

\section{Holding}

Expanding \textit{Bivens} is a disfavored judicial activity. A court first asks whether a claim arises in a new context---whether it differs in a meaningful way from \textit{Bivens}, \textit{Davis v. Passman}, or \textit{Carlson v. Green}. If so, no remedy may be implied when special factors counsel hesitation. The detention-policy claims were barred; the abuse claim against Warden Hasty arose in a new context and was remanded for special-factors analysis; and all petitioners had qualified immunity on the \S~1985(3) claim.

\section{Analysis of Opinions}

\textbf{Majority (Kennedy).} \textit{Bivens}, \textit{Davis}, and \textit{Carlson} were products of an ``ancien regime'' that freely implied causes of action. Preserving \textit{Bivens} in its core sphere, Kennedy held that Congress, not the courts, should usually decide whether to create damages liability. Indicators of a new context include the officers' rank, the right at issue, the generality of official action, judicial guidance, statutory mandate, and the risk of intrusion on other branches. Four special factors barred the policy claims: \textit{Bivens} is not a vehicle to alter high-level policy, and discovery would chill deliberation; the claims implicate national security; Congress's long silence despite oversight and an Inspector General report suggested deliberate refusal; and injunctions and habeas were potential alternatives. On Hasty, ``even a modest extension is still an extension'': the claim rested on the Fifth Amendment, and the Prison Litigation Reform Act omitted a damages remedy. Officials could not have known that intra-departmental consultation was an unlawful conspiracy.

\textbf{Thomas (concurring in part).} \textit{Bivens} is a ``relic'' that should be confined to its facts; he also urged reconsidering modern qualified immunity.

\textbf{Breyer (dissenting).} The claims arose in familiar \textit{Bivens} and \textit{Carlson} territory; Congress ratified \textit{Bivens} in the Westfall Act; blackouts left detainees with damages or nothing; and \textit{Iqbal}, qualified immunity, and tailored discovery already protect officials in emergencies.

\textbf{Later.} \textit{Hernandez v. Mesa} (2020), \textit{Egbert v. Boule} (2022), and \textit{Goldey v. Fields} (2025) tightened the framework further.

\section{Examples: Future Applications}

\begin{enumerate}
\item \textbf{Same-Side (No Remedy): Emergency Directive.} Travelers sue the DHS Secretary over warrantless device sweeps ordered after a cyber-threat---high-level policy and national security.
\item \textbf{Same-Side (No Remedy): Retaliation.} A motel owner sues an agent for retaliatory investigations; the Court has never recognized a First Amendment \textit{Bivens} claim.
\item \textbf{Opposite-Side (Likely Available): Warrantless Entry.} An ATF agent breaks down a door without a warrant and punches the homeowner---close to \textit{Bivens} itself.
\item \textbf{Opposite-Side (Likely Available): Prison Medical Care.} Officials deny an asthmatic prisoner care until he dies, tracking \textit{Carlson}.
\item \textbf{Fence-Sitter: Pretrial Detainee.} A Marshal beats a detainee and denies insulin; the medical claim tests whether the Fifth/Eighth Amendment difference alone makes a new context.
\end{enumerate}

\section{Critique}

\textbf{Rights without remedies.} State officers face damages under \S~1983 for the same violations from which federal officers increasingly escape, and treating injunctions or grievance procedures as alternatives ignores that many detainees face damages or nothing.

\textbf{Overruling in slow motion.} Because the new-context inquiry is easily satisfied and any plausible concern counsels hesitation, critics argue that \textit{Abbasi} and \textit{Egbert} confine \textit{Bivens} to its 1971 facts without candidly overruling it.

\section{Key Quotations}

\begin{quote}
``The proper test for determining whether a case presents a new \textit{Bivens} context is as follows. If the case is different in a meaningful way from previous \textit{Bivens} cases decided by this Court, then the context is new.''
\end{quote}
\noindent --- Justice Kennedy, for the Court, 582 U.S. at 139.

\begin{quote}
``Given the notable change in the Court's approach to recognizing implied causes of action, however, the Court has made clear that expanding the \textit{Bivens} remedy is now a `disfavored' judicial activity.''
\end{quote}
\noindent --- Justice Kennedy, for the Court, 582 U.S. at 135.

\begin{quote}
``The necessary inference, though, is that the inquiry must concentrate on whether the Judiciary is well suited, absent congressional action or instruction, to consider and weigh the costs and benefits of allowing a damages action to proceed. Thus, to be a `special factor counselling hesitation,' a factor must cause a court to hesitate before answering that question in the affirmative.''
\end{quote}
\noindent --- Justice Kennedy, for the Court, 582 U.S. at 136.

\end{document}
